% --- Archon physics formalization source begin ---
% archon:physics
% archon:covers PhyXMiniProblems/problem_phyx_mini_0581.lean
% archon:source-report reports/phyx_mini/problem_phyx_mini_0581.source.json

\chapter{Physics problem phyx\_mini\_0581}
\label{ch:PhyXMiniProblems_problem_phyx_mini_0581}

\paragraph{Problem source.}
\#\# Physical scenario

Figure is an energy-level diagram for a fictitious infinite potential well that contains one electron. The number of degenerate states of the levels are indicated: "non" means nondegenerate (which includes the ground state of the electron), "double" means 2 states, and "triple" means 3 states. We put a total of 11 electrons in the well.

\#\# Question

If the electrostatic forces between the electrons can be neglected, what multiple of \textbackslash{}( h\textasciicircum{}2 \textbackslash{}!/ 8mL\textasciicircum{}2 \textbackslash{}) gives the energy of the first excited state of the 11-electron system?

\#\# Answer choices

- A: A: 65 \textbackslash{}left( \textbackslash{}frac\{h\textasciicircum{}2\}\{8mL\textasciicircum{}2\} 
ight)

- B: B: 68 \textbackslash{}left( \textbackslash{}frac\{h\textasciicircum{}2\}\{8mL\textasciicircum{}2\} 
ight)

- C: C: 66 \textbackslash{}left( \textbackslash{}frac\{h\textasciicircum{}2\}\{8mL\textasciicircum{}2\} 
ight)

- D: D: 64 \textbackslash{}left( \textbackslash{}frac\{h\textasciicircum{}2\}\{8mL\textasciicircum{}2\} 
ight)

\#\# Figure caption (auxiliary; use the image as primary evidence)

The image shows a diagram representing energy levels labeled "E (h²/8mL²)" on the left side. There are horizontal lines indicating these energy levels at specific values, with labels on the right:

- A line at level "12" labeled "Non."
- A line at level "11" labeled "Triple."
- Lines at levels "7" and "6," both labeled "Double" and "Triple."
- A line at level "4" labeled "Ground."

These lines represent quantized energy states, typical in quantum mechanics, where the numerical values correspond to discrete energy levels.

\#\# Dataset metadata

Quantum Phenomena; Physical Model Grounding Reasoning, Multi-Formula Reasoning

\paragraph{Recorded answer/context.}
C: C: 66 \textbackslash{}left( \textbackslash{}frac\{h\textasciicircum{}2\}\{8mL\textasciicircum{}2\} 
ight)

\paragraph{Figure/image path.}
/root/proposal\_for\_physic/science-mango/phyx\_mini\_run/phyx\_data/test\_image/581.png

\paragraph{Formalization target.}
create a compiling Lean file with sorry bodies at `PhyXMiniProblems/problem\_phyx\_mini\_0581.lean`. The Lean declarations must preserve the physical quantities, dimensions or dimensional roles, figure labels, governing-law hypotheses, and final relation expressed by this problem.
Use Mathlib/Physlib names found through LeanExplore where available. If a physics API is missing, introduce faithful local abstractions rather than scalar placeholder aliases.

\begin{theorem}[Physics formalization target]
\label{thm:physics:phyx_mini_0581:target}
\lean{PhyXMiniProblems.ProblemPhyXMini0581.firstExcitedEnergyOfElevenElectrons_is_choiceC}
\uses{def:physics:phyx-mini-0581:phyxminiproblems-problemphyxmini0581-isenergymultipleof, def:physics:phyx-mini-0581:phyxminiproblems-problemphyxmini0581-elevenelectronwellsetup, def:physics:phyx-mini-0581:phyxminiproblems-problemphyxmini0581-matcheselevenelectronwellscenario, def:physics:phyx-mini-0581:phyxminiproblems-problemphyxmini0581-matchessuppliedinfinitewelldiagram, def:physics:phyx-mini-0581:phyxminiproblems-problemphyxmini0581-satisfiesinfinitewellenergyscalelaw, def:physics:phyx-mini-0581:phyxminiproblems-problemphyxmini0581-usesordinaryplanckconstant, def:physics:phyx-mini-0581:phyxminiproblems-problemphyxmini0581-hasphysicalinfinitewellparameters, def:physics:phyx-mini-0581:phyxminiproblems-problemphyxmini0581-satisfiesnoninteractingelectronenergylaw, def:physics:phyx-mini-0581:phyxminiproblems-problemphyxmini0581-hasgroundandfirstexcitedconfigurations, def:physics:phyx-mini-0581:phyxminiproblems-problemphyxmini0581-answerchoice, def:physics:phyx-mini-0581:phyxminiproblems-problemphyxmini0581-answerchoice-energymultiplier, def:physics:phyx-mini-0581:phyxminiproblems-problemphyxmini0581-matchesdisplayedfirstexcitedenergy}
The assigned autoformalize agent should translate this physics problem into Lean declarations in the covered file, with theorem and lemma proof bodies written as `by sorry`.
\end{theorem}
\begin{proof}
This is an autoformalization task, not a proof task. Produce statements that can later be proved by the physics prover without weakening the physical model.
\end{proof}

% --- Archon declaration topology begin ---
\section{Declaration topology}
\begin{definition}[Length Quantity]
  \label{def:physics:phyx-mini-0581:phyxminiproblems-problemphyxmini0581-lengthquantity}
  \lean{PhyXMiniProblems.ProblemPhyXMini0581.LengthQuantity}
  A nonnegative, unit-independent physical length
\end{definition}
\begin{proof}
  This is the declared model definition of Length Quantity.
\end{proof}
\begin{definition}[Mass Quantity]
  \label{def:physics:phyx-mini-0581:phyxminiproblems-problemphyxmini0581-massquantity}
  \lean{PhyXMiniProblems.ProblemPhyXMini0581.MassQuantity}
  A nonnegative, unit-independent physical mass
\end{definition}
\begin{proof}
  This is the declared model definition of Mass Quantity.
\end{proof}
\begin{definition}[action Dimension]
  \label{def:physics:phyx-mini-0581:phyxminiproblems-problemphyxmini0581-actiondimension}
  \lean{PhyXMiniProblems.ProblemPhyXMini0581.actionDimension}
  The physical dimension of action, used for the ordinary Planck constant
\end{definition}
\begin{proof}
  This is the declared model definition of action Dimension.
\end{proof}
\begin{definition}[Action Quantity]
  \label{def:physics:phyx-mini-0581:phyxminiproblems-problemphyxmini0581-actionquantity}
  \lean{PhyXMiniProblems.ProblemPhyXMini0581.ActionQuantity}
  \uses{def:physics:phyx-mini-0581:phyxminiproblems-problemphyxmini0581-actiondimension}
  A nonnegative, unit-independent physical action
\end{definition}
\begin{proof}
  This is the declared model definition of Action Quantity.
\end{proof}
\begin{definition}[length Readout]
  \label{def:physics:phyx-mini-0581:phyxminiproblems-problemphyxmini0581-lengthreadout}
  \lean{PhyXMiniProblems.ProblemPhyXMini0581.lengthReadout}
  \uses{def:physics:phyx-mini-0581:phyxminiproblems-problemphyxmini0581-lengthquantity}
  Read a physical length in the selected length unit
\end{definition}
\begin{proof}
  This is the declared model definition of length Readout.
\end{proof}
\begin{definition}[mass Readout]
  \label{def:physics:phyx-mini-0581:phyxminiproblems-problemphyxmini0581-massreadout}
  \lean{PhyXMiniProblems.ProblemPhyXMini0581.massReadout}
  \uses{def:physics:phyx-mini-0581:phyxminiproblems-problemphyxmini0581-massquantity}
  Read a physical mass in the selected mass unit
\end{definition}
\begin{proof}
  This is the declared model definition of mass Readout.
\end{proof}
\begin{definition}[length In Meters]
  \label{def:physics:phyx-mini-0581:phyxminiproblems-problemphyxmini0581-lengthinmeters}
  \lean{PhyXMiniProblems.ProblemPhyXMini0581.lengthInMeters}
  \uses{def:physics:phyx-mini-0581:phyxminiproblems-problemphyxmini0581-lengthquantity, def:physics:phyx-mini-0581:phyxminiproblems-problemphyxmini0581-lengthreadout}
  Width of the well, read in metres
\end{definition}
\begin{proof}
  This is the declared model definition of length In Meters.
\end{proof}
\begin{definition}[mass In Kilograms]
  \label{def:physics:phyx-mini-0581:phyxminiproblems-problemphyxmini0581-massinkilograms}
  \lean{PhyXMiniProblems.ProblemPhyXMini0581.massInKilograms}
  \uses{def:physics:phyx-mini-0581:phyxminiproblems-problemphyxmini0581-massquantity, def:physics:phyx-mini-0581:phyxminiproblems-problemphyxmini0581-massreadout}
  Electron mass, read in kilograms
\end{definition}
\begin{proof}
  This is the declared model definition of mass In Kilograms.
\end{proof}
\begin{definition}[action In Joule Seconds]
  \label{def:physics:phyx-mini-0581:phyxminiproblems-problemphyxmini0581-actioninjouleseconds}
  \lean{PhyXMiniProblems.ProblemPhyXMini0581.actionInJouleSeconds}
  \uses{def:physics:phyx-mini-0581:phyxminiproblems-problemphyxmini0581-actionquantity}
  Ordinary Planck action, read coherently in joule-seconds
\end{definition}
\begin{proof}
  This is the declared model definition of action In Joule Seconds.
\end{proof}
\begin{definition}[energy Readout]
  \label{def:physics:phyx-mini-0581:phyxminiproblems-problemphyxmini0581-energyreadout}
  \lean{PhyXMiniProblems.ProblemPhyXMini0581.energyReadout}
  Read an energy in an arbitrary coherent choice of units
\end{definition}
\begin{proof}
  This is the declared model definition of energy Readout.
\end{proof}
\begin{definition}[energy In Joules]
  \label{def:physics:phyx-mini-0581:phyxminiproblems-problemphyxmini0581-energyinjoules}
  \lean{PhyXMiniProblems.ProblemPhyXMini0581.energyInJoules}
  \uses{def:physics:phyx-mini-0581:phyxminiproblems-problemphyxmini0581-energyreadout}
  Read an energy in coherent SI units, hence in joules
\end{definition}
\begin{proof}
  This is the declared model definition of energy In Joules.
\end{proof}
\begin{definition}[Is Energy Multiple Of]
  \label{def:physics:phyx-mini-0581:phyxminiproblems-problemphyxmini0581-isenergymultipleof}
  \lean{PhyXMiniProblems.ProblemPhyXMini0581.IsEnergyMultipleOf}
  \uses{def:physics:phyx-mini-0581:phyxminiproblems-problemphyxmini0581-energyreadout}
  energy is multiplier times scale independently of the units used to read the two energies. This is the dimensionally meaningful form of the multiple asked for in the problem
\end{definition}
\begin{proof}
  This is the declared model definition of Is Energy Multiple Of.
\end{proof}
\begin{definition}[Well Energy Level]
  \label{def:physics:phyx-mini-0581:phyxminiproblems-problemphyxmini0581-wellenergylevel}
  \lean{PhyXMiniProblems.ProblemPhyXMini0581.WellEnergyLevel}
  The five one-electron levels shown in the supplied diagram
\end{definition}
\begin{proof}
  This is the declared model definition of Well Energy Level.
\end{proof}
\begin{definition}[Printed Level Annotation]
  \label{def:physics:phyx-mini-0581:phyxminiproblems-problemphyxmini0581-printedlevelannotation}
  \lean{PhyXMiniProblems.ProblemPhyXMini0581.PrintedLevelAnnotation}
  Literal kinds of annotations printed to the right of the level lines
\end{definition}
\begin{proof}
  This is the declared model definition of Printed Level Annotation.
\end{proof}
\begin{definition}[spatial Multiplicity]
  \label{def:physics:phyx-mini-0581:phyxminiproblems-problemphyxmini0581-printedlevelannotation-spatialmultiplicity}
  \lean{PhyXMiniProblems.ProblemPhyXMini0581.PrintedLevelAnnotation.spatialMultiplicity}
  \uses{def:physics:phyx-mini-0581:phyxminiproblems-problemphyxmini0581-printedlevelannotation}
  Orbital multiplicity represented by each printed annotation
\end{definition}
\begin{proof}
  This is the declared model definition of spatial Multiplicity.
\end{proof}
\begin{definition}[Infinite Well Energy Diagram]
  \label{def:physics:phyx-mini-0581:phyxminiproblems-problemphyxmini0581-infinitewellenergydiagram}
  \lean{PhyXMiniProblems.ProblemPhyXMini0581.InfiniteWellEnergyDiagram}
  \uses{def:physics:phyx-mini-0581:phyxminiproblems-problemphyxmini0581-wellenergylevel, def:physics:phyx-mini-0581:phyxminiproblems-problemphyxmini0581-printedlevelannotation}
  Primary-image data. The axis label records that the displayed vertical numbers are multiples of h\textasciicircum{}2 / (8mL\textasciicircum{}2) rather than bare energies
\end{definition}
\begin{proof}
  This is the declared model definition of Infinite Well Energy Diagram.
\end{proof}
\begin{definition}[Particle Species]
  \label{def:physics:phyx-mini-0581:phyxminiproblems-problemphyxmini0581-particlespecies}
  \lean{PhyXMiniProblems.ProblemPhyXMini0581.ParticleSpecies}
  Particle species represented by the occupation numbers
\end{definition}
\begin{proof}
  This is the declared model definition of Particle Species.
\end{proof}
\begin{definition}[Potential Model]
  \label{def:physics:phyx-mini-0581:phyxminiproblems-problemphyxmini0581-potentialmodel}
  \lean{PhyXMiniProblems.ProblemPhyXMini0581.PotentialModel}
  Potential model selected in the fictitious problem
\end{definition}
\begin{proof}
  This is the declared model definition of Potential Model.
\end{proof}
\begin{definition}[Electron Configuration]
  \label{def:physics:phyx-mini-0581:phyxminiproblems-problemphyxmini0581-electronconfiguration}
  \lean{PhyXMiniProblems.ProblemPhyXMini0581.ElectronConfiguration}
  \uses{def:physics:phyx-mini-0581:phyxminiproblems-problemphyxmini0581-wellenergylevel}
  Number of electrons occupying each displayed one-particle energy level
\end{definition}
\begin{proof}
  This is the declared model definition of Electron Configuration.
\end{proof}
\begin{definition}[Eleven Electron Well Setup]
  \label{def:physics:phyx-mini-0581:phyxminiproblems-problemphyxmini0581-elevenelectronwellsetup}
  \lean{PhyXMiniProblems.ProblemPhyXMini0581.ElevenElectronWellSetup}
  \uses{def:physics:phyx-mini-0581:phyxminiproblems-problemphyxmini0581-lengthquantity, def:physics:phyx-mini-0581:phyxminiproblems-problemphyxmini0581-massquantity, def:physics:phyx-mini-0581:phyxminiproblems-problemphyxmini0581-actionquantity, def:physics:phyx-mini-0581:phyxminiproblems-problemphyxmini0581-wellenergylevel, def:physics:phyx-mini-0581:phyxminiproblems-problemphyxmini0581-infinitewellenergydiagram, def:physics:phyx-mini-0581:phyxminiproblems-problemphyxmini0581-particlespecies, def:physics:phyx-mini-0581:phyxminiproblems-problemphyxmini0581-potentialmodel, def:physics:phyx-mini-0581:phyxminiproblems-problemphyxmini0581-electronconfiguration}
  Independent physical quantities and states of the model. Neither the ground nor first-excited configuration is defined by an answer choice or by a numerical occupation pattern; the predicates below characterize them by Pauli admissibility and energy order
\end{definition}
\begin{proof}
  This is the declared model definition of Eleven Electron Well Setup.
\end{proof}
\begin{definition}[configuration Electron Count]
  \label{def:physics:phyx-mini-0581:phyxminiproblems-problemphyxmini0581-configurationelectroncount}
  \lean{PhyXMiniProblems.ProblemPhyXMini0581.configurationElectronCount}
  \uses{def:physics:phyx-mini-0581:phyxminiproblems-problemphyxmini0581-wellenergylevel, def:physics:phyx-mini-0581:phyxminiproblems-problemphyxmini0581-electronconfiguration}
  Total number of electrons in an occupation-number configuration
\end{definition}
\begin{proof}
  This is the declared model definition of configuration Electron Count.
\end{proof}
\begin{definition}[pauli Capacity]
  \label{def:physics:phyx-mini-0581:phyxminiproblems-problemphyxmini0581-paulicapacity}
  \lean{PhyXMiniProblems.ProblemPhyXMini0581.pauliCapacity}
  \uses{def:physics:phyx-mini-0581:phyxminiproblems-problemphyxmini0581-wellenergylevel, def:physics:phyx-mini-0581:phyxminiproblems-problemphyxmini0581-elevenelectronwellsetup}
  Pauli capacity of a level: spin multiplicity times orbital degeneracy
\end{definition}
\begin{proof}
  This is the declared model definition of pauli Capacity.
\end{proof}
\begin{definition}[Is Admissible Configuration]
  \label{def:physics:phyx-mini-0581:phyxminiproblems-problemphyxmini0581-isadmissibleconfiguration}
  \lean{PhyXMiniProblems.ProblemPhyXMini0581.IsAdmissibleConfiguration}
  \uses{def:physics:phyx-mini-0581:phyxminiproblems-problemphyxmini0581-electronconfiguration, def:physics:phyx-mini-0581:phyxminiproblems-problemphyxmini0581-elevenelectronwellsetup, def:physics:phyx-mini-0581:phyxminiproblems-problemphyxmini0581-configurationelectroncount, def:physics:phyx-mini-0581:phyxminiproblems-problemphyxmini0581-paulicapacity}
  An eleven-electron configuration obeys particle-number conservation and the Pauli capacity of every displayed level
\end{definition}
\begin{proof}
  This is the declared model definition of Is Admissible Configuration.
\end{proof}
\begin{definition}[configuration Energy In Joules]
  \label{def:physics:phyx-mini-0581:phyxminiproblems-problemphyxmini0581-configurationenergyinjoules}
  \lean{PhyXMiniProblems.ProblemPhyXMini0581.configurationEnergyInJoules}
  \uses{def:physics:phyx-mini-0581:phyxminiproblems-problemphyxmini0581-energyinjoules, def:physics:phyx-mini-0581:phyxminiproblems-problemphyxmini0581-electronconfiguration, def:physics:phyx-mini-0581:phyxminiproblems-problemphyxmini0581-elevenelectronwellsetup}
  SI energy readout of a many-electron occupation configuration
\end{definition}
\begin{proof}
  This is the declared model definition of configuration Energy In Joules.
\end{proof}
\begin{definition}[Is Ground Configuration]
  \label{def:physics:phyx-mini-0581:phyxminiproblems-problemphyxmini0581-isgroundconfiguration}
  \lean{PhyXMiniProblems.ProblemPhyXMini0581.IsGroundConfiguration}
  \uses{def:physics:phyx-mini-0581:phyxminiproblems-problemphyxmini0581-electronconfiguration, def:physics:phyx-mini-0581:phyxminiproblems-problemphyxmini0581-elevenelectronwellsetup, def:physics:phyx-mini-0581:phyxminiproblems-problemphyxmini0581-isadmissibleconfiguration, def:physics:phyx-mini-0581:phyxminiproblems-problemphyxmini0581-configurationenergyinjoules}
  An admissible configuration of least many-electron energy
\end{definition}
\begin{proof}
  This is the declared model definition of Is Ground Configuration.
\end{proof}
\begin{definition}[Is First Excited Configuration]
  \label{def:physics:phyx-mini-0581:phyxminiproblems-problemphyxmini0581-isfirstexcitedconfiguration}
  \lean{PhyXMiniProblems.ProblemPhyXMini0581.IsFirstExcitedConfiguration}
  \uses{def:physics:phyx-mini-0581:phyxminiproblems-problemphyxmini0581-electronconfiguration, def:physics:phyx-mini-0581:phyxminiproblems-problemphyxmini0581-elevenelectronwellsetup, def:physics:phyx-mini-0581:phyxminiproblems-problemphyxmini0581-isadmissibleconfiguration, def:physics:phyx-mini-0581:phyxminiproblems-problemphyxmini0581-configurationenergyinjoules, def:physics:phyx-mini-0581:phyxminiproblems-problemphyxmini0581-isgroundconfiguration}
  An admissible state of least energy strictly above a specified ground state. This definition permits degeneracy at the first-excited energy
\end{definition}
\begin{proof}
  This is the declared model definition of Is First Excited Configuration.
\end{proof}
\begin{definition}[Matches Eleven Electron Well Scenario]
  \label{def:physics:phyx-mini-0581:phyxminiproblems-problemphyxmini0581-matcheselevenelectronwellscenario}
  \lean{PhyXMiniProblems.ProblemPhyXMini0581.MatchesElevenElectronWellScenario}
  \uses{def:physics:phyx-mini-0581:phyxminiproblems-problemphyxmini0581-elevenelectronwellsetup}
  Prose data and modeling regime stated in the problem
\end{definition}
\begin{proof}
  This is the declared model definition of Matches Eleven Electron Well Scenario.
\end{proof}
\begin{definition}[Matches Supplied Infinite Well Diagram]
  \label{def:physics:phyx-mini-0581:phyxminiproblems-problemphyxmini0581-matchessuppliedinfinitewelldiagram}
  \lean{PhyXMiniProblems.ProblemPhyXMini0581.MatchesSuppliedInfiniteWellDiagram}
  \uses{def:physics:phyx-mini-0581:phyxminiproblems-problemphyxmini0581-energyreadout, def:physics:phyx-mini-0581:phyxminiproblems-problemphyxmini0581-elevenelectronwellsetup}
  Exact level values and annotations transcribed from the primary image, together with the calibration of physical one-particle energies and degeneracies from those readouts. No many-electron energy or occupation pattern occurs here
\end{definition}
\begin{proof}
  This is the declared model definition of Matches Supplied Infinite Well Diagram.
\end{proof}
\begin{definition}[Satisfies Infinite Well Energy Scale Law]
  \label{def:physics:phyx-mini-0581:phyxminiproblems-problemphyxmini0581-satisfiesinfinitewellenergyscalelaw}
  \lean{PhyXMiniProblems.ProblemPhyXMini0581.SatisfiesInfiniteWellEnergyScaleLaw}
  \uses{def:physics:phyx-mini-0581:phyxminiproblems-problemphyxmini0581-lengthinmeters, def:physics:phyx-mini-0581:phyxminiproblems-problemphyxmini0581-massinkilograms, def:physics:phyx-mini-0581:phyxminiproblems-problemphyxmini0581-actioninjouleseconds, def:physics:phyx-mini-0581:phyxminiproblems-problemphyxmini0581-energyinjoules, def:physics:phyx-mini-0581:phyxminiproblems-problemphyxmini0581-elevenelectronwellsetup}
  The displayed energy unit is the physical quantity h²/(8mL²). The equality is expressed using coherent SI readouts, where its right side is in joules. This scale law does not specify any many-electron state or answer multiplier
\end{definition}
\begin{proof}
  This is the declared model definition of Satisfies Infinite Well Energy Scale Law.
\end{proof}
\begin{definition}[Uses Ordinary Planck Constant]
  \label{def:physics:phyx-mini-0581:phyxminiproblems-problemphyxmini0581-usesordinaryplanckconstant}
  \lean{PhyXMiniProblems.ProblemPhyXMini0581.UsesOrdinaryPlanckConstant}
  \uses{def:physics:phyx-mini-0581:phyxminiproblems-problemphyxmini0581-actioninjouleseconds, def:physics:phyx-mini-0581:phyxminiproblems-problemphyxmini0581-elevenelectronwellsetup}
  Physlib exposes the reduced Planck constant Constants.ℏ; the ordinary constant in the figure is h = 2πℏ. This reference fact remains symbolic with respect to the requested integer multiplier
\end{definition}
\begin{proof}
  This is the declared model definition of Uses Ordinary Planck Constant.
\end{proof}
\begin{definition}[Has Physical Infinite Well Parameters]
  \label{def:physics:phyx-mini-0581:phyxminiproblems-problemphyxmini0581-hasphysicalinfinitewellparameters}
  \lean{PhyXMiniProblems.ProblemPhyXMini0581.HasPhysicalInfiniteWellParameters}
  \uses{def:physics:phyx-mini-0581:phyxminiproblems-problemphyxmini0581-lengthinmeters, def:physics:phyx-mini-0581:phyxminiproblems-problemphyxmini0581-massinkilograms, def:physics:phyx-mini-0581:phyxminiproblems-problemphyxmini0581-actioninjouleseconds, def:physics:phyx-mini-0581:phyxminiproblems-problemphyxmini0581-energyinjoules, def:physics:phyx-mini-0581:phyxminiproblems-problemphyxmini0581-elevenelectronwellsetup}
  Positivity assumptions selecting a physically meaningful well
\end{definition}
\begin{proof}
  This is the declared model definition of Has Physical Infinite Well Parameters.
\end{proof}
\begin{definition}[Satisfies Noninteracting Electron Energy Law]
  \label{def:physics:phyx-mini-0581:phyxminiproblems-problemphyxmini0581-satisfiesnoninteractingelectronenergylaw}
  \lean{PhyXMiniProblems.ProblemPhyXMini0581.SatisfiesNoninteractingElectronEnergyLaw}
  \uses{def:physics:phyx-mini-0581:phyxminiproblems-problemphyxmini0581-energyreadout, def:physics:phyx-mini-0581:phyxminiproblems-problemphyxmini0581-wellenergylevel, def:physics:phyx-mini-0581:phyxminiproblems-problemphyxmini0581-electronconfiguration, def:physics:phyx-mini-0581:phyxminiproblems-problemphyxmini0581-elevenelectronwellsetup, def:physics:phyx-mini-0581:phyxminiproblems-problemphyxmini0581-isadmissibleconfiguration}
  When electrostatic interactions are neglected, the many-electron energy is the sum of occupied one-electron energies. This law is stated for every admissible configuration and in every unit choice; it contains no special occupation pattern and no answer multiplier
\end{definition}
\begin{proof}
  This is the declared model definition of Satisfies Noninteracting Electron Energy Law.
\end{proof}
\begin{definition}[Has Ground And First Excited Configurations]
  \label{def:physics:phyx-mini-0581:phyxminiproblems-problemphyxmini0581-hasgroundandfirstexcitedconfigurations}
  \lean{PhyXMiniProblems.ProblemPhyXMini0581.HasGroundAndFirstExcitedConfigurations}
  \uses{def:physics:phyx-mini-0581:phyxminiproblems-problemphyxmini0581-elevenelectronwellsetup, def:physics:phyx-mini-0581:phyxminiproblems-problemphyxmini0581-isgroundconfiguration, def:physics:phyx-mini-0581:phyxminiproblems-problemphyxmini0581-isfirstexcitedconfiguration}
  The two distinguished configuration fields really denote the ground and first-excited states according to the general energy-order definitions above. No numerical energy or occupation is assumed
\end{definition}
\begin{proof}
  This is the declared model definition of Has Ground And First Excited Configurations.
\end{proof}
\begin{definition}[Answer Choice]
  \label{def:physics:phyx-mini-0581:phyxminiproblems-problemphyxmini0581-answerchoice}
  \lean{PhyXMiniProblems.ProblemPhyXMini0581.AnswerChoice}
  Labels of the four displayed energy-multiplier choices
\end{definition}
\begin{proof}
  This is the declared model definition of Answer Choice.
\end{proof}
\begin{definition}[energy Multiplier]
  \label{def:physics:phyx-mini-0581:phyxminiproblems-problemphyxmini0581-answerchoice-energymultiplier}
  \lean{PhyXMiniProblems.ProblemPhyXMini0581.AnswerChoice.energyMultiplier}
  \uses{def:physics:phyx-mini-0581:phyxminiproblems-problemphyxmini0581-answerchoice}
  Integer multiple of h²/(8mL²) printed beside an answer choice
\end{definition}
\begin{proof}
  This is the declared model definition of energy Multiplier.
\end{proof}
\begin{definition}[recorded Dataset Answer]
  \label{def:physics:phyx-mini-0581:phyxminiproblems-problemphyxmini0581-recordeddatasetanswer}
  \lean{PhyXMiniProblems.ProblemPhyXMini0581.recordedDatasetAnswer}
  \uses{def:physics:phyx-mini-0581:phyxminiproblems-problemphyxmini0581-answerchoice}
  Answer label recorded in the source dataset
\end{definition}
\begin{proof}
  This is the declared model definition of recorded Dataset Answer.
\end{proof}
\begin{definition}[Matches Displayed First Excited Energy]
  \label{def:physics:phyx-mini-0581:phyxminiproblems-problemphyxmini0581-matchesdisplayedfirstexcitedenergy}
  \lean{PhyXMiniProblems.ProblemPhyXMini0581.MatchesDisplayedFirstExcitedEnergy}
  \uses{def:physics:phyx-mini-0581:phyxminiproblems-problemphyxmini0581-isenergymultipleof, def:physics:phyx-mini-0581:phyxminiproblems-problemphyxmini0581-elevenelectronwellsetup, def:physics:phyx-mini-0581:phyxminiproblems-problemphyxmini0581-answerchoice}
  A displayed choice matches the physical first-excited system energy
\end{definition}
\begin{proof}
  This is the declared model definition of Matches Displayed First Excited Energy.
\end{proof}
\begin{lemma}[ground configuration occupation and energy]
  \label{lem:physics:phyx-mini-0581:phyxminiproblems-problemphyxmini0581-ground-configuration-occupation-and-energy}
  \lean{PhyXMiniProblems.ProblemPhyXMini0581.ground_configuration_occupation_and_energy}
  \uses{def:physics:phyx-mini-0581:phyxminiproblems-problemphyxmini0581-isenergymultipleof, def:physics:phyx-mini-0581:phyxminiproblems-problemphyxmini0581-elevenelectronwellsetup, def:physics:phyx-mini-0581:phyxminiproblems-problemphyxmini0581-matcheselevenelectronwellscenario, def:physics:phyx-mini-0581:phyxminiproblems-problemphyxmini0581-matchessuppliedinfinitewelldiagram, def:physics:phyx-mini-0581:phyxminiproblems-problemphyxmini0581-hasphysicalinfinitewellparameters, def:physics:phyx-mini-0581:phyxminiproblems-problemphyxmini0581-satisfiesnoninteractingelectronenergylaw, def:physics:phyx-mini-0581:phyxminiproblems-problemphyxmini0581-hasgroundandfirstexcitedconfigurations}
  Filling the lowest Pauli capacities gives the eleven-electron ground occupation (2, 6, 3, 0, 0) and total multiplier 2·4 + 6·6 + 3·7 = 65
\end{lemma}
\begin{proof}
  The conclusion follows from the governing relations and figure readouts named in the statement of ground configuration occupation and energy.
\end{proof}
\begin{lemma}[first excited configuration occupation]
  \label{lem:physics:phyx-mini-0581:phyxminiproblems-problemphyxmini0581-first-excited-configuration-occupation}
  \lean{PhyXMiniProblems.ProblemPhyXMini0581.first_excited_configuration_occupation}
  \uses{def:physics:phyx-mini-0581:phyxminiproblems-problemphyxmini0581-elevenelectronwellsetup, def:physics:phyx-mini-0581:phyxminiproblems-problemphyxmini0581-matcheselevenelectronwellscenario, def:physics:phyx-mini-0581:phyxminiproblems-problemphyxmini0581-matchessuppliedinfinitewelldiagram, def:physics:phyx-mini-0581:phyxminiproblems-problemphyxmini0581-hasphysicalinfinitewellparameters, def:physics:phyx-mini-0581:phyxminiproblems-problemphyxmini0581-satisfiesnoninteractingelectronenergylaw, def:physics:phyx-mini-0581:phyxminiproblems-problemphyxmini0581-hasgroundandfirstexcitedconfigurations}
  The least energy strictly above the ground state moves one electron from the filled multiplier-6 level into the one remaining spin-orbital at multiplier 7, giving occupation (2, 5, 4, 0, 0)
\end{lemma}
\begin{proof}
  The conclusion follows from the governing relations and figure readouts named in the statement of first excited configuration occupation.
\end{proof}
% --- Archon declaration topology end ---
% --- Archon physics formalization source end ---
